% ============================================================
% Energy-Consistency Analysis Section Draft
% ============================================================
% This file contains only the analysis section intended to be
% inserted into a larger manuscript.
%
% Required theorem-like environments in the manuscript preamble:
%   \newtheorem{assumption}{Assumption}
%   \newtheorem{lemma}{Lemma}
%   \newtheorem{proposition}{Proposition}
%   \newtheorem{theorem}{Theorem}
%   \newtheorem{remark}{Remark}
%
% Required packages typically include:
%   amsmath, amssymb, amsthm, mathtools
%
% Notation warning:
%   We write \partial_x w and \partial_y w for derivatives.
%   The symbols u_x and u_y denote represented solutions, not
%   partial derivatives of u.
% ============================================================

\documentclass[11pt]{article}

\usepackage[margin=1in]{geometry}
\usepackage{amsmath,amssymb,amsthm,mathtools}
\usepackage{enumitem}
\usepackage{bm}

\newtheorem{assumption}{Assumption}
\newtheorem{lemma}{Lemma}
\newtheorem{proposition}{Proposition}
\newtheorem{theorem}{Theorem}
\newtheorem{remark}{Remark}

\begin{document}

\section{Energy-Consistency Analysis for Linewise Green Reconstructions}
\label{sec:energy-consistency-analysis}

This section analyzes how an energy-consistency loss between two
directional reconstructions can control the source-split error variables
used in the proposed linewise Green reconstruction framework.  The
analysis is carried out in a continuous setting.  The main purpose is not
to claim that the learned Green operators are exact, but to isolate the
mathematical mechanism by which the consistency loss
\[
    \mathcal L_E
    =
    \int_\Omega a |\nabla(u_x-u_y)|^2
\]
controls the average and directional components of the split error under
explicit structural assumptions.  We also discuss why an \(L^2\)-based
consistency loss does not yield the same argument, and how the conclusion
is perturbed when the Green reconstructions are not exact.

\subsection{Setting and directional operators}
\label{subsec:setting-directional-operators}

Let
\[
    \Omega=(0,1)^2,
    \qquad
    V=H_0^1(\Omega).
\]
We consider a directional decomposition of a second-order operator
\[
    L = L_x + L_y,
\]
where
\begin{align}
    L_x u
    &:=
    -\partial_x\!\left(a\,\partial_x u\right)
    + b_x\,\partial_x u
    + c u,
    \label{eq:Lx-definition}
    \\
    L_y u
    &:=
    -\partial_y\!\left(a\,\partial_y u\right)
    + b_y\,\partial_y u
    + c u.
    \label{eq:Ly-definition}
\end{align}
Thus the full operator is
\begin{equation}
    L u
    =
    -\nabla\cdot(a\nabla u)
    + b\cdot\nabla u
    + 2c u,
    \qquad
    b=(b_x,b_y).
    \label{eq:L-full-definition}
\end{equation}
Let \(u_*\in V\) denote the reference solution satisfying
\begin{equation}
    L u_* = f
    \label{eq:reference-problem}
\end{equation}
in the appropriate weak sense.

Throughout this section, the energy seminorm induced by the diffusion
coefficient is denoted by
\begin{equation}
    \|v\|_a^2
    :=
    \int_\Omega a |\nabla v|^2.
    \label{eq:energy-norm}
\end{equation}
Under the ellipticity assumption \(a\ge a_0>0\), this seminorm is
equivalent to the \(H_0^1\)-seminorm on \(V\).

\subsection{Fiberwise one-dimensional Green reconstructions}
\label{subsec:fiberwise-green-reconstructions}

The operators \(G_x\) and \(G_y\) used below should not be understood as
two-dimensional inverses of the full operator \(L\).  Instead, they are
fiberwise one-dimensional Green inverses acting on two-dimensional source
fields.

For a two-dimensional source field \(\phi=\phi(x,y)\), define
\[
    \phi_y(\xi) := \phi(\xi,y).
\]
For almost every \(y\in(0,1)\), let \(G_x^{(y)}\) denote the
one-dimensional Green inverse associated with the \(x\)-directional
operator
\[
    L_x^{(y)} w
    =
    -\partial_x\!\left(a(x,y)\partial_x w\right)
    + b_x(x,y)\partial_x w
    + c(x,y)w
\]
on the interval \((0,1)\), with homogeneous endpoint conditions
\[
    w(0)=w(1)=0.
\]
Then the fiberwise reconstruction \(G_x[\phi]\) is defined by
\begin{equation}
    (G_x\phi)(x,y)
    :=
    \bigl(G_x^{(y)}\phi_y\bigr)(x)
    \qquad
    \text{for a.e. } y\in(0,1).
    \label{eq:Gx-fiberwise-definition}
\end{equation}
Equivalently, when the one-dimensional Green kernel exists,
\[
    (G_x\phi)(x,y)
    =
    \int_0^1 G_x(x,\xi;y)\,\phi(\xi,y)\,d\xi .
\]

Similarly, for \(\psi=\psi(x,y)\), define
\[
    \psi_x(\eta) := \psi(x,\eta).
\]
For almost every \(x\in(0,1)\), let \(G_y^{(x)}\) denote the
one-dimensional Green inverse associated with
\[
    L_y^{(x)} w
    =
    -\partial_y\!\left(a(x,y)\partial_y w\right)
    + b_y(x,y)\partial_y w
    + c(x,y)w
\]
on \((0,1)\), with homogeneous endpoint conditions
\[
    w(0)=w(1)=0.
\]
Then
\begin{equation}
    (G_y\psi)(x,y)
    :=
    \bigl(G_y^{(x)}\psi_x\bigr)(y)
    \qquad
    \text{for a.e. } x\in(0,1).
    \label{eq:Gy-fiberwise-definition}
\end{equation}
Equivalently,
\[
    (G_y\psi)(x,y)
    =
    \int_0^1 G_y(y,\eta;x)\,\psi(x,\eta)\,d\eta .
\]

\begin{remark}[Fiberwise reconstruction versus full two-dimensional admissibility]
The reconstruction \(G_x[\phi]\) naturally enforces homogeneous boundary
conditions only on the \(x\)-endpoints of each fiber:
\[
    (G_x\phi)(0,y)=(G_x\phi)(1,y)=0
    \qquad \text{for a.e. }y.
\]
Likewise, \(G_y[\psi]\) naturally enforces homogeneous boundary
conditions only on the \(y\)-endpoints:
\[
    (G_y\psi)(x,0)=(G_y\psi)(x,1)=0
    \qquad \text{for a.e. }x.
\]
Thus, from the fiberwise construction alone, one should generally expect
mapping properties of the form
\[
    G_x[\phi]\in L^2_y(0,1;H_0^1(0,1)_x),
    \qquad
    G_y[\psi]\in L^2_x(0,1;H_0^1(0,1)_y),
\]
rather than the full two-dimensional admissibility
\[
    G_x[\phi],\,G_y[\psi]\in H_0^1(\Omega).
\]
The latter property will therefore be imposed explicitly as a structural
admissibility assumption in the theorem below.
\end{remark}

\subsection{Exact and predicted source splits}
\label{subsec:exact-predicted-splits}

The reference solution \(u_*\) induces an exact directional source split
\begin{equation}
    \phi_* := L_x u_*,
    \qquad
    \psi_* := L_y u_*.
    \label{eq:exact-split-definition}
\end{equation}
Consequently,
\begin{equation}
    \phi_*+\psi_*
    =
    (L_x+L_y)u_*
    =
    Lu_*
    =
    f.
    \label{eq:exact-split-consistency}
\end{equation}
The model predicts another split \((\phi,\psi)\), and the corresponding
represented solutions are defined by
\begin{equation}
    u_x := G_x[\phi],
    \qquad
    u_y := G_y[\psi].
    \label{eq:represented-solutions}
\end{equation}
The model architecture is intended to enforce the split consistency
constraint
\begin{equation}
    \phi+\psi=f,
    \label{eq:predicted-split-consistency}
\end{equation}
up to numerical error.  In the continuous analysis below, we impose this
constraint exactly.

\subsection{Structural assumptions}
\label{subsec:structural-assumptions}

We now state the assumptions used in the energy-consistency analysis.

\begin{assumption}[Coefficient bounds]
\label{assumption:coefficient-bounds}
There exist constants \(a_0,a_1>0\) such that
\begin{equation}
    0<a_0\le a(x,y)\le a_1<\infty
    \qquad
    \text{for a.e. }(x,y)\in\Omega.
    \label{eq:a-bounds}
\end{equation}
Moreover,
\[
    b\in L^\infty(\Omega)^2,
    \qquad
    c\in L^\infty(\Omega).
\]
Unless otherwise stated, \(\|b\|_{L^\infty}\) denotes the vector
\(L^\infty\)-norm
\[
    \|b\|_{L^\infty}
    :=
    \operatorname*{ess\,sup}_{(x,y)\in\Omega}
    \sqrt{|b_x(x,y)|^2+|b_y(x,y)|^2}.
\]
\end{assumption}

\begin{assumption}[Boundary admissibility]
\label{assumption:boundary-admissibility}
The reference and represented solutions satisfy
\begin{equation}
    u_x,u_y,u_*\in H_0^1(\Omega).
    \label{eq:boundary-admissibility}
\end{equation}
This assumption is automatic for \(u_*\) when the reference problem is a
homogeneous Dirichlet problem.  For \(u_x\) and \(u_y\), however, it is
not a consequence of the fiberwise Green reconstruction alone; it should
be interpreted as a restriction to a two-dimensional admissible class.
\end{assumption}

\begin{assumption}[Exact Green reconstruction of the reference split]
\label{assumption:exact-green-reconstruction}
The fiberwise Green inverses exactly reconstruct the reference solution
from the exact split:
\begin{equation}
    G_x[\phi_*]=u_*,
    \qquad
    G_y[\psi_*]=u_*.
    \label{eq:exact-green-reconstruction}
\end{equation}
This assumption corresponds to an ideal Green reconstruction setting.
A perturbation version with imperfect Green reconstruction is discussed
in Subsection~\ref{subsec:imperfect-green-reconstruction}.
\end{assumption}

\begin{assumption}[Projection consistency]
\label{assumption:projection-consistency}
The predicted split satisfies
\begin{equation}
    \phi+\psi=f=\phi_*+\psi_*.
    \label{eq:projection-consistency}
\end{equation}
Equivalently,
\begin{equation}
    (\phi-\phi_*)+(\psi-\psi_*)=0.
    \label{eq:projection-consistency-error}
\end{equation}
\end{assumption}

\begin{assumption}[Weak inverse identities]
\label{assumption:weak-inverse-identities}
The fiberwise Green reconstructions satisfy the directional inverse
identities in the weak sense.  In particular, for appropriate source
terms \(\mu,\nu\),
\begin{equation}
    L_x G_x[\mu]=\mu,
    \qquad
    L_y G_y[\nu]=\nu.
    \label{eq:weak-inverse-identities}
\end{equation}
Equivalently, in terms of the directional bilinear forms defined below,
\[
    B_x(G_x[\mu],v)=\langle \mu,v\rangle,
    \qquad
    B_y(G_y[\nu],v)=\langle \nu,v\rangle
\]
for all admissible test functions \(v\).  In particular,
\begin{equation}
    L_x q_x = 2(\phi-\phi_*),
    \qquad
    L_y q_y = 2(\psi-\psi_*),
    \label{eq:inverse-identities-for-q}
\end{equation}
for the error variables \(q_x,q_y\) defined in
Subsection~\ref{subsec:error-variables}.
\end{assumption}

\begin{assumption}[Coercivity of the full bilinear form]
\label{assumption:coercivity}
Let
\begin{equation}
    B(u,v)
    :=
    \int_\Omega a\nabla u\cdot\nabla v
    +
    \int_\Omega (b\cdot\nabla u)v
    +
    \int_\Omega 2cuv.
    \label{eq:full-bilinear-form}
\end{equation}
There exists a constant \(\gamma_a>0\) such that
\begin{equation}
    B(v,v)
    \ge
    \gamma_a \|v\|_a^2
    \qquad
    \forall v\in V.
    \label{eq:coercivity}
\end{equation}
\end{assumption}

\subsection{Coercivity of the full bilinear form}
\label{subsec:coercivity-discussion}

We briefly discuss when Assumption~\ref{assumption:coercivity} is
reasonable.  For \(v\in V\),
\begin{equation}
    B(v,v)
    =
    \|v\|_a^2
    +
    \int_\Omega (b\cdot\nabla v)v
    +
    \int_\Omega 2cv^2.
    \label{eq:Bvv-expanded}
\end{equation}
Thus coercivity means that the lower-order advection and reaction terms
do not destroy the positive diffusion energy.

In the pure diffusion case, \(b=0\) and \(c=0\), one has
\[
    B(v,v)=\|v\|_a^2,
\]
so \(\gamma_a=1\).  Similarly, if \(b=0\) and \(c\ge0\), then
\[
    B(v,v)=\|v\|_a^2+\int_\Omega 2cv^2
    \ge
    \|v\|_a^2,
\]
again with \(\gamma_a=1\).

If the reaction coefficient can be negative, write
\[
    c_-(x,y):=\max\{-c(x,y),0\}.
\]
Then
\[
    \int_\Omega 2cv^2
    \ge
    -2\|c_-\|_{L^\infty}\|v\|_{L^2}^2.
\]
Using Poincaré's inequality
\[
    \|v\|_{L^2}\le C_P\|\nabla v\|_{L^2}
\]
and the ellipticity bound \(a\ge a_0\), we get
\[
    \|v\|_{L^2}^2
    \le
    \frac{C_P^2}{a_0}\|v\|_a^2.
\]
Therefore, in the reaction-diffusion case \(b=0\),
\[
    B(v,v)
    \ge
    \left(
    1-\frac{2C_P^2\|c_-\|_{L^\infty}}{a_0}
    \right)
    \|v\|_a^2.
\]
Hence a sufficient condition for coercivity is
\[
    2C_P^2\|c_-\|_{L^\infty}<a_0.
\]

When \(b\neq0\), a conservative sufficient condition can be obtained by
bounding the advection term as
\begin{align*}
    \left|
    \int_\Omega (b\cdot\nabla v)v
    \right|
    &\le
    \|b\|_{L^\infty}
    \|\nabla v\|_{L^2}
    \|v\|_{L^2}
    \\
    &\le
    \frac{C_P\|b\|_{L^\infty}}{a_0}
    \|v\|_a^2.
\end{align*}
Combining this with the negative part of the reaction term yields
\[
    B(v,v)
    \ge
    \left(
    1
    -
    \frac{C_P\|b\|_{L^\infty}}{a_0}
    -
    \frac{2C_P^2\|c_-\|_{L^\infty}}{a_0}
    \right)
    \|v\|_a^2.
\]
Thus one sufficient, although generally conservative, condition is
\begin{equation}
    C_P\|b\|_{L^\infty}
    +
    2C_P^2\|c_-\|_{L^\infty}
    <
    a_0.
    \label{eq:conservative-coercivity-condition}
\end{equation}

A sharper interpretation is possible when \(b\) is sufficiently regular,
for example \(b\in W^{1,\infty}(\Omega)^2\).  Since \(v\in H_0^1(\Omega)\),
integration by parts gives
\begin{equation}
    \int_\Omega (b\cdot\nabla v)v
    =
    \frac12\int_\Omega b\cdot\nabla(v^2)
    =
    -\frac12\int_\Omega (\nabla\cdot b)v^2.
    \label{eq:advection-integration-by-parts}
\end{equation}
Hence
\begin{equation}
    B(v,v)
    =
    \|v\|_a^2
    +
    \int_\Omega
    \left(
    2c-\frac12\nabla\cdot b
    \right)v^2.
    \label{eq:Bvv-effective-reaction}
\end{equation}
If
\[
    2c-\frac12\nabla\cdot b\ge0
    \qquad
    \text{a.e. in }\Omega,
\]
then \(B(v,v)\ge \|v\|_a^2\), and \(\gamma_a=1\).  More generally, if
\[
    2c-\frac12\nabla\cdot b\ge -\kappa
\]
with
\[
    \kappa < \frac{a_0}{C_P^2},
\]
then
\[
    B(v,v)
    \ge
    \left(
    1-\frac{\kappa C_P^2}{a_0}
    \right)
    \|v\|_a^2.
\]
In particular, if \(\nabla\cdot b=0\) and \(c\ge0\), the advection term
does not contribute to \(B(v,v)\), and coercivity follows with
\(\gamma_a=1\).

\begin{remark}[Role of coercivity]
Coercivity is the step that turns the weak identity for \(q_c\) into an
energy estimate.  Without a lower bound of the form
\[
    B(q_c,q_c)\ge \gamma_a\|q_c\|_a^2,
\]
the boundedness of the right-hand side alone would not control
\(\|q_c\|_a\).
\end{remark}

\subsection{Error variables and average--difference decomposition}
\label{subsec:error-variables}

Define the directional split-error variables
\begin{equation}
    q_x:=2G_x(\phi-\phi_*),
    \qquad
    q_y:=2G_y(\psi-\psi_*).
    \label{eq:qx-qy-definition}
\end{equation}
Also define the average component and the difference component by
\begin{equation}
    q_c:=\frac{q_x+q_y}{2},
    \qquad
    r:=u_x-u_y.
    \label{eq:qc-r-definition}
\end{equation}
By linearity of the Green reconstructions and by
Assumption~\ref{assumption:exact-green-reconstruction},
\begin{align}
    q_x
    &=
    2\bigl(G_x[\phi]-G_x[\phi_*]\bigr)
    =
    2(u_x-u_*),
    \label{eq:qx-solution-error}
    \\
    q_y
    &=
    2\bigl(G_y[\psi]-G_y[\psi_*]\bigr)
    =
    2(u_y-u_*).
    \label{eq:qy-solution-error}
\end{align}
Therefore,
\[
    q_x-q_y
    =
    2(u_x-u_y)
    =
    2r.
\]
Consequently,
\begin{equation}
    q_x=q_c+r,
    \qquad
    q_y=q_c-r.
    \label{eq:average-difference-decomposition}
\end{equation}
The energy-consistency loss is precisely
\begin{equation}
    \mathcal L_E
    :=
    \int_\Omega a|\nabla(u_x-u_y)|^2
    =
    \|r\|_a^2.
    \label{eq:energy-loss-r}
\end{equation}

\subsection{A structural identity for the average component}
\label{subsec:structural-identity}

The next lemma identifies the equation satisfied by the average component
\(q_c\).

\begin{lemma}[Structural identity for the average component]
\label{lemma:structural-identity}
Under Assumptions~\ref{assumption:projection-consistency} and
\ref{assumption:weak-inverse-identities}, the average component \(q_c\)
satisfies
\begin{equation}
    Lq_c=-(L_x-L_y)r
    \label{eq:structural-identity}
\end{equation}
in the weak sense.
\end{lemma}

\begin{proof}
By the weak inverse identities,
\[
    L_xq_x=2(\phi-\phi_*),
    \qquad
    L_yq_y=2(\psi-\psi_*).
\]
Adding these two equations gives
\[
    L_xq_x+L_yq_y
    =
    2\{(\phi-\phi_*)+(\psi-\psi_*)\}.
\]
By projection consistency,
\[
    (\phi-\phi_*)+(\psi-\psi_*)=0.
\]
Thus
\[
    L_xq_x+L_yq_y=0.
\]
Using the decomposition
\[
    q_x=q_c+r,
    \qquad
    q_y=q_c-r,
\]
we obtain
\[
    L_x(q_c+r)+L_y(q_c-r)=0.
\]
Equivalently,
\[
    (L_x+L_y)q_c+(L_x-L_y)r=0.
\]
Since \(L=L_x+L_y\), this gives
\[
    Lq_c=-(L_x-L_y)r.
\]
\end{proof}

\subsection{Boundedness of the directional difference}
\label{subsec:boundedness-directional-difference}

Define the directional bilinear forms
\begin{align}
    B_x(w,v)
    & :=
    \int_\Omega a\,\partial_x w\,\partial_x v
    +
    \int_\Omega b_x\,\partial_x w\,v
    +
    \int_\Omega cwv,
    \label{eq:Bx-definition}
    \\
    B_y(w,v)
    & :=
    \int_\Omega a\,\partial_y w\,\partial_y v
    +
    \int_\Omega b_y\,\partial_y w\,v
    +
    \int_\Omega cwv.
    \label{eq:By-definition}
\end{align}
Then
\[
    B(w,v)=B_x(w,v)+B_y(w,v).
\]

\begin{lemma}[Boundedness of the directional difference]
\label{lemma:boundedness-directional-difference}
For all \(r,v\in H_0^1(\Omega)\),
\begin{equation}
    |B_x(r,v)-B_y(r,v)|
    \le
    M\|r\|_a\|v\|_a,
    \label{eq:directional-difference-bound}
\end{equation}
where
\begin{equation}
    M
    =
    1+
    \frac{C_P\|b\|_{L^\infty}}{a_0}.
    \label{eq:M-constant}
\end{equation}
Here \(C_P\) is the Poincaré constant for \(H_0^1(\Omega)\):
\[
    \|v\|_{L^2}\le C_P\|\nabla v\|_{L^2}.
\]
\end{lemma}

\begin{proof}
By definition,
\begin{align}
    B_x(r,v)-B_y(r,v)
    &=
    \int_\Omega a
    \bigl(\partial_x r\,\partial_x v
          -\partial_y r\,\partial_y v\bigr)
    \notag\\
    &\quad
    +
    \int_\Omega
    \bigl(b_x\partial_x r-b_y\partial_y r\bigr)v.
    \label{eq:Bx-minus-By-expanded}
\end{align}
The reaction terms cancel because the same term \(\int_\Omega crv\)
appears in both \(B_x(r,v)\) and \(B_y(r,v)\).

We first estimate the diffusion part.  Write
\[
    P_r:=(\partial_x r,-\partial_y r),
    \qquad
    P_v:=(\partial_x v,\partial_y v).
\]
Then
\[
    P_r\cdot P_v
    =
    \partial_x r\,\partial_x v
    -
    \partial_y r\,\partial_y v,
\]
and
\[
    |P_r|=|\nabla r|,
    \qquad
    |P_v|=|\nabla v|.
\]
Therefore, by the weighted Cauchy--Schwarz inequality,
\begin{align}
    \left|
    \int_\Omega a
    \bigl(\partial_x r\,\partial_x v
          -\partial_y r\,\partial_y v\bigr)
    \right|
    &=
    \left|
    \int_\Omega a\,P_r\cdot P_v
    \right|
    \notag\\
    &\le
    \left(\int_\Omega a|P_r|^2\right)^{1/2}
    \left(\int_\Omega a|P_v|^2\right)^{1/2}
    \notag\\
    &=
    \|r\|_a\|v\|_a.
    \label{eq:diffusion-part-bound}
\end{align}

Next, define \(\widetilde b=(b_x,-b_y)\).  Since
\(|\widetilde b|=|b|\), we have
\[
    \|\widetilde b\|_{L^\infty}=\|b\|_{L^\infty}.
\]
The advection part can be written as
\[
    \int_\Omega
    \bigl(b_x\partial_x r-b_y\partial_y r\bigr)v
    =
    \int_\Omega
    (\widetilde b\cdot\nabla r)v.
\]
Using Hölder's inequality and then Cauchy--Schwarz,
\begin{align}
    \left|
    \int_\Omega
    (\widetilde b\cdot\nabla r)v
    \right|
    &\le
    \|b\|_{L^\infty}
    \|\nabla r\|_{L^2}
    \|v\|_{L^2}.
    \label{eq:advection-preliminary-bound}
\end{align}
By \(a\ge a_0>0\),
\[
    \|\nabla r\|_{L^2}
    \le
    \frac{1}{\sqrt{a_0}}\|r\|_a.
\]
Moreover, by Poincaré's inequality and \(a\ge a_0\),
\[
    \|v\|_{L^2}
    \le
    C_P\|\nabla v\|_{L^2}
    \le
    \frac{C_P}{\sqrt{a_0}}\|v\|_a.
\]
Substituting these estimates into
\eqref{eq:advection-preliminary-bound} gives
\begin{equation}
    \left|
    \int_\Omega
    \bigl(b_x\partial_x r-b_y\partial_y r\bigr)v
    \right|
    \le
    \frac{C_P\|b\|_{L^\infty}}{a_0}
    \|r\|_a\|v\|_a.
    \label{eq:advection-part-bound}
\end{equation}
Combining \eqref{eq:diffusion-part-bound} and
\eqref{eq:advection-part-bound}, we obtain
\[
    |B_x(r,v)-B_y(r,v)|
    \le
    \left(
    1+
    \frac{C_P\|b\|_{L^\infty}}{a_0}
    \right)
    \|r\|_a\|v\|_a.
\]
This proves the claim.
\end{proof}

\subsection{Main energy-consistency theorem}
\label{subsec:main-energy-consistency-theorem}

We now combine the structural identity, coercivity, and the boundedness
of the directional difference.

\begin{theorem}[Energy-consistency control of the split errors]
\label{theorem:energy-consistency-control}
Suppose Assumptions~\ref{assumption:coefficient-bounds}--
\ref{assumption:coercivity} hold.  Let
\[
    \mathcal L_E
    =
    \|u_x-u_y\|_a^2
    =
    \|r\|_a^2.
\]
Then
\begin{equation}
    \|q_c\|_a
    \le
    C_E\sqrt{\mathcal L_E},
    \label{eq:qc-energy-control}
\end{equation}
where
\begin{equation}
    C_E
    =
    \frac{M}{\gamma_a}
    =
    \frac{
    1+\dfrac{C_P\|b\|_{L^\infty}}{a_0}
    }{\gamma_a}.
    \label{eq:CE-definition}
\end{equation}
Furthermore,
\begin{equation}
    \|q_x\|_a
    \le
    (1+C_E)\sqrt{\mathcal L_E},
    \qquad
    \|q_y\|_a
    \le
    (1+C_E)\sqrt{\mathcal L_E}.
    \label{eq:qx-qy-energy-control}
\end{equation}
Consequently,
\begin{equation}
    \mathcal L_E\to0
    \quad\Longrightarrow\quad
    q_c,q_x,q_y\to0
    \quad
    \text{in the energy norm}.
    \label{eq:energy-loss-convergence}
\end{equation}
Since \(q_c,q_x,q_y\in H_0^1(\Omega)\), energy-norm convergence also
implies convergence in \(H_0^1(\Omega)\).
\end{theorem}

\begin{proof}
By Lemma~\ref{lemma:structural-identity},
\[
    Lq_c=-(L_x-L_y)r
\]
in the weak sense.  Therefore,
\begin{equation}
    B(q_c,v)
    =
    -\{B_x(r,v)-B_y(r,v)\}
    \qquad
    \forall v\in V.
    \label{eq:weak-form-structural-identity}
\end{equation}
Taking \(v=q_c\) gives
\[
    B(q_c,q_c)
    =
    -\{B_x(r,q_c)-B_y(r,q_c)\}.
\]
By coercivity,
\[
    \gamma_a\|q_c\|_a^2
    \le
    B(q_c,q_c).
\]
Hence
\[
    \gamma_a\|q_c\|_a^2
    \le
    |B_x(r,q_c)-B_y(r,q_c)|.
\]
By Lemma~\ref{lemma:boundedness-directional-difference},
\[
    |B_x(r,q_c)-B_y(r,q_c)|
    \le
    M\|r\|_a\|q_c\|_a.
\]
Thus
\[
    \gamma_a\|q_c\|_a^2
    \le
    M\|r\|_a\|q_c\|_a.
\]
If \(\|q_c\|_a=0\), then the desired estimate is trivial.  Otherwise,
dividing by \(\|q_c\|_a\) yields
\[
    \|q_c\|_a
    \le
    \frac{M}{\gamma_a}\|r\|_a.
\]
Since \(\mathcal L_E=\|r\|_a^2\), we obtain
\[
    \|q_c\|_a
    \le
    C_E\sqrt{\mathcal L_E}.
\]

It remains to estimate \(q_x\) and \(q_y\).  From
\[
    q_x=q_c+r,
    \qquad
    q_y=q_c-r,
\]
the triangle inequality gives
\[
    \|q_x\|_a
    \le
    \|q_c\|_a+\|r\|_a
    \le
    (C_E+1)\sqrt{\mathcal L_E},
\]
and similarly
\[
    \|q_y\|_a
    \le
    \|q_c\|_a+\|r\|_a
    \le
    (C_E+1)\sqrt{\mathcal L_E}.
\]
This proves the stated estimates.

Finally, if \(\mathcal L_E\to0\), then
\[
    \sqrt{\mathcal L_E}\to0,
\]
and the above estimates imply
\[
    q_c,q_x,q_y\to0
    \quad
    \text{in the energy norm}.
\]
Because \(a\ge a_0>0\), energy-norm convergence implies convergence of
the gradients in \(L^2\).  Since the functions lie in \(H_0^1(\Omega)\),
Poincaré's inequality also controls their \(L^2\)-norms.  Hence the
convergence also holds in \(H_0^1(\Omega)\).
\end{proof}

\subsection{Why an \(L^2\)-consistency loss is insufficient}
\label{subsec:L2-consistency-insufficient}

The preceding argument relies essentially on the energy norm.  Suppose
instead that one used the \(L^2\)-consistency loss
\[
    \mathcal L_0
    :=
    \|u_x-u_y\|_{L^2(\Omega)}^2
    =
    \|r\|_{L^2(\Omega)}^2.
\]
This loss directly controls only the amplitude of \(r\), not its
derivatives.  However, the right-hand side of the weak structural
identity involves
\[
    B_x(r,v)-B_y(r,v),
\]
which contains derivatives of \(r\):
\begin{align*}
    B_x(r,v)-B_y(r,v)
    &=
    \int_\Omega a
    \bigl(\partial_x r\,\partial_x v
          -\partial_y r\,\partial_y v\bigr)
    \\
    &\quad
    +
    \int_\Omega
    \bigl(b_x\partial_x r-b_y\partial_y r\bigr)v.
\end{align*}
Therefore, the current proof requires control of \(\nabla r\), naturally
provided by \(\|r\|_a\), but not by \(\|r\|_{L^2}\).

In general, one cannot replace Lemma~\ref{lemma:boundedness-directional-difference}
by an estimate of the form
\[
    |B_x(r,v)-B_y(r,v)|
    \le
    C\|r\|_{L^2}\|v\|_a.
\]
A simple high-frequency example illustrates the obstruction.  In the
constant-coefficient pure diffusion case, let
\[
    r_n(x,y)=\sin(n\pi x)\sin(\pi y).
\]
Then
\[
    \|r_n\|_{L^2}=O(1),
    \qquad
    \|\partial_x r_n\|_{L^2}=O(n).
\]
Thus the derivative-dependent terms in \(B_x(r_n,v)-B_y(r_n,v)\) can
grow while the \(L^2\)-norm of \(r_n\) remains bounded.  Consequently,
\(L^2\)-consistency alone does not control \((L_x-L_y)r\), and hence does
not yield the same estimate for \(q_c\) through the present weak-form
argument.

\begin{remark}
This does not rule out all possible \(L^2\)-based estimates under
additional structure.  For example, in very special constant-coefficient
settings, one may attempt spectral or elliptic-regularity arguments for
the operator \(L^{-1}(L_x-L_y)\).  Such an argument is different from the
energy method used here and would require additional regularity and
boundary analysis.  The present theorem explains why the energy norm is
the natural norm for the current proof.
\end{remark}

\subsection{Perturbation by imperfect Green reconstruction}
\label{subsec:imperfect-green-reconstruction}

The main theorem assumes exact Green reconstruction of the reference
split:
\[
    G_x[\phi_*]=u_*,
    \qquad
    G_y[\psi_*]=u_*.
\]
In practice, however, the Green operators may be learned or approximated.
We therefore consider the perturbed setting
\begin{equation}
    G_x[\phi_*]=u_*+\varepsilon_x,
    \qquad
    G_y[\psi_*]=u_*+\varepsilon_y,
    \label{eq:imperfect-green-reconstruction}
\end{equation}
where \(\varepsilon_x,\varepsilon_y\) measure the reference
reconstruction errors.

Assuming linearity of the approximate reconstructions, the same formal
definitions
\[
    q_x:=2G_x(\phi-\phi_*),
    \qquad
    q_y:=2G_y(\psi-\psi_*)
\]
now give
\begin{align}
    q_x
    &=
    2(G_x[\phi]-G_x[\phi_*])
    =
    2(u_x-u_*-\varepsilon_x),
    \label{eq:qx-imperfect-green}
    \\
    q_y
    &=
    2(G_y[\psi]-G_y[\psi_*])
    =
    2(u_y-u_*-\varepsilon_y).
    \label{eq:qy-imperfect-green}
\end{align}
Define the branch mismatch error
\begin{equation}
    \delta\varepsilon
    :=
    \varepsilon_x-\varepsilon_y.
    \label{eq:branch-mismatch-error}
\end{equation}
Then
\[
    q_x-q_y
    =
    2(u_x-u_y)-2(\varepsilon_x-\varepsilon_y)
    =
    2(r-\delta\varepsilon).
\]
Thus the effective difference component is
\begin{equation}
    s:=r-\delta\varepsilon.
    \label{eq:effective-difference-component}
\end{equation}
With
\[
    q_c:=\frac{q_x+q_y}{2},
\]
we obtain the perturbed decomposition
\begin{equation}
    q_x=q_c+s,
    \qquad
    q_y=q_c-s.
    \label{eq:perturbed-average-difference-decomposition}
\end{equation}

If the weak inverse identities remain valid for the error variables
\(q_x,q_y\), then the same projection argument yields
\[
    L_xq_x+L_yq_y=0.
\]
Substituting \eqref{eq:perturbed-average-difference-decomposition} gives
\[
    L_x(q_c+s)+L_y(q_c-s)=0,
\]
and therefore
\begin{equation}
    Lq_c=-(L_x-L_y)s.
    \label{eq:perturbed-structural-identity}
\end{equation}
Applying the same coercivity and boundedness argument as in
Theorem~\ref{theorem:energy-consistency-control}, we obtain
\[
    \|q_c\|_a
    \le
    C_E\|s\|_a.
\]
Since
\[
    \|s\|_a
    =
    \|r-\delta\varepsilon\|_a
    \le
    \|r\|_a+\|\delta\varepsilon\|_a,
\]
and \(\|r\|_a=\sqrt{\mathcal L_E}\), we get
\begin{equation}
    \|q_c\|_a
    \le
    C_E
    \left(
    \sqrt{\mathcal L_E}
    +
    \|\varepsilon_x-\varepsilon_y\|_a
    \right).
    \label{eq:perturbed-qc-estimate}
\end{equation}
Similarly,
\begin{equation}
    \|q_x\|_a,\ \|q_y\|_a
    \le
    (1+C_E)
    \left(
    \sqrt{\mathcal L_E}
    +
    \|\varepsilon_x-\varepsilon_y\|_a
    \right).
    \label{eq:perturbed-qx-qy-estimate}
\end{equation}

However, the actual represented-solution errors include the individual
Green reconstruction errors:
\begin{equation}
    2(u_x-u_*)=q_x+2\varepsilon_x,
    \qquad
    2(u_y-u_*)=q_y+2\varepsilon_y.
    \label{eq:solution-error-with-green-bias}
\end{equation}
Thus, even if \(\mathcal L_E\) is small, the represented solution errors
may not be small unless the Green reconstruction errors are also
controlled.

\begin{remark}[Common Green reconstruction bias]
Energy consistency penalizes branch disagreement, not common Green
reconstruction bias.  If
\[
    \varepsilon_x=\varepsilon_y=\varepsilon,
\]
then
\[
    \varepsilon_x-\varepsilon_y=0,
\]
so the branch mismatch term does not appear in
\eqref{eq:perturbed-qc-estimate}.  Nevertheless,
\[
    2(u_x-u_*)=q_x+2\varepsilon,
    \qquad
    2(u_y-u_*)=q_y+2\varepsilon.
\]
Therefore, a common Green reconstruction bias can remain invisible to
the consistency loss \( \|u_x-u_y\|_a^2 \), while still contributing to
the actual solution error.
\end{remark}

\begin{remark}[Possible inverse residuals for learned Green operators]
The perturbation discussion above accounts for reference reconstruction
errors but assumes that the weak inverse identities remain valid for the
error variables.  A learned Green operator may also introduce inverse
residuals.  For example, if
\[
    L_x\widehat G_x[\mu]
    =
    \mu+\eta_x(\mu),
    \qquad
    L_y\widehat G_y[\nu]
    =
    \nu+\eta_y(\nu),
\]
then the structural identity acquires an additional residual term.  In a
dual energy norm, this would lead to an additional contribution of the
form
\[
    \frac{1}{\gamma_a}\|R\|_{a,*},
\]
where \(R\) collects the directional inverse residuals and
\[
    \|R\|_{a,*}
    :=
    \sup_{v\in V\setminus\{0\}}
    \frac{|\langle R,v\rangle|}{\|v\|_a}.
\]
Thus, in a fully learned setting, the consistency estimate should be
interpreted as controlling the split errors up to Green reconstruction
error and inverse residual error.
\end{remark}

\subsection{Summary of the analysis}
\label{subsec:analysis-summary}

The analysis shows that, under the stated admissibility, projection,
Green reconstruction, weak inverse, and coercivity assumptions, the
energy-consistency loss controls the split-error variables:
\[
    \mathcal L_E\to0
    \quad\Longrightarrow\quad
    q_c,q_x,q_y\to0
    \quad
    \text{in the energy norm}.
\]
The mechanism is as follows.  The loss directly controls the difference
component \(r=u_x-u_y\).  The projection constraint and weak inverse
identities yield the structural identity
\[
    Lq_c=-(L_x-L_y)r,
\]
which allows the average component \(q_c\) to be controlled by \(r\)
through coercivity of \(L\) and boundedness of \(B_x-B_y\).  Finally, the
directional split errors \(q_x,q_y\) are controlled by the
average--difference decomposition
\[
    q_x=q_c+r,
    \qquad
    q_y=q_c-r.
\]
The same argument also clarifies the limitations: an \(L^2\)-consistency
loss does not control the derivative terms in \(B_x-B_y\), and imperfect
Green reconstruction introduces an error floor that must be separately
controlled.

\end{document}
